\section{Discussion}
\label{sec:discussion}

The results support the two-level account and, with it, the gate view of replay. This section reads them back against the theory, weighs the limits of the evidence, and draws out what follows for method design.

\subsection{Why the gap needs two levels}
\label{subsec:disc_twolevel}

The central empirical claim is that the stability gap is not one quantity but two failures of the same optimisation, kept apart because they answer to different interventions, and the $2\times2$ of Section~\ref{subsec:res_levels} is the direct test: removing the overshoot drivers leaves the arc, fixing the path leaves the overshoot, and only doing both closes the gap. Neither intervention is a weaker version of the other. The arc is a property of which path the descent dynamics define, and no amount of optimiser hygiene (exact gradients, balanced magnitudes, zero noise) touches it; the overshoot is how far the discrete, noisy, accelerated iterate leaves whatever path it is on, and no change of path removes it. An intervention aimed at one level is silent on the other.

The natural alternative, a single additive decomposition of the gap into magnitude, variance, and trajectory shares, does not survive the data: the shares are order-dependent (a noisy replay gradient inflates the first-step magnitude), and the mapping from factors to metrics is crossed by the experiments, the curriculum's momentum-$0$ tail looking like a ``trajectory'' term until C8 removes it with the exact buffer and no momentum at all, revealing it as variance (Table~\ref{tab:res_c8}). The two-level split is instead by kind: which path is followed versus how tightly it is followed.

The split also settles what the gap is not the price of: C8 closes the transient to $0.2$ pp while its accuracy sits $5.5$ pp below vanilla ER (Section~\ref{subsec:res_levels}), so the permanent rise to the joint-basin level is the stability--plasticity cost and everything below the settled level and back is failure to follow the monotone reference path, not the cost of learning. This is why the end-referenced area is anchored to the recovered plateau rather than the pre-switch baseline: it attributes the gap to the trajectory rather than to capacity loss.

Finally, the two levels explain why the metric that matters changes with the backbone. On the MLP the network reconverges in a few steps, so the arc integrates to little and the perceived severity is the overshoot spike; on the ResNet the same transient stays open for hundreds of steps (Figure~\ref{fig:cifar_curves}), so the slow tail is where most of the lost accuracy accumulates and where both gates' margin over vanilla is largest (Table~\ref{tab:res_cifar}). Depth and area are the signatures of the two levels, and a depth-only reading would have understated the result on exactly the architectures where continual-evaluation cost is highest.

\subsection{Two gates, two ways to divide the work}
\label{subsec:disc_gates}

Written in the unified update of Section~\ref{subsec:gates}, every replay method gates the new-task push while restoring at full strength, and the curriculum and asymmetric PER are the two pure corners of the taxonomy: a scalar gate driven by feedback, and a directional gate driven by feedforward. The results show the two corners working and failing exactly where the taxonomy says they should, and the single fact that organises the whole comparison is \emph{where} the gate attenuates the push.

The scalar feedback gate attenuates at source: it multiplies the entire new-task loss by a schedule that starts near zero, so early in the transition there is almost no push to discretise, accumulate, or let noise act on. Three consequences follow, and the data confirm each. It composes with momentum, because the accumulator never sees a sustained full-strength push to rebuild (Section~\ref{subsec:res_scalar}). It has no curvature blind spot, because the adaptive schedule reads a live damage signal rather than a model of the past task. And it pays a plasticity debt: holding the loss down slows the new task, which is why C8 at momentum $0$ is gap-free but $5.5$ pp down on accuracy, a debt only acceleration repays (momentum lifts C8 from $0.818$ to $0.955$, Table~\ref{tab:res_c8}). The scalar gate's price is plasticity, and momentum is how it is paid.

The directional feedforward gate attenuates per step: it filters the push through the replay Fisher, never slows the new task at source, and so protects the past task at no plasticity cost, ACC and short-horizon plasticity rising as the filter strengthens (Table~\ref{tab:res_per}). The same per-step locality is its weakness on two axes, one the taxonomy predicts and one that emerged unanticipated. Under momentum the protection unravels by the re-inflation mechanism of Section~\ref{subsec:res_directional}, depth climbing from ${\approx}2$ to ${\approx}8$ pp: the same lever that closes the scalar gate's gap re-opens the directional gate's, and the difference is entirely where the attenuation is applied, at the source or at the step. That one distinction accounts for the whole momentum column of the study: the vanilla null, the curriculum's composition, and the directional gate's re-inflation.

The directional gate's second failure was not predicted, and it is the sharper one because it makes the case for feedback. At the strongest filter the past-task curve can slide off a late cliff (the $\delta{=}0.03$ full-buffer cliffs of Section~\ref{subsec:res_directional}). A plausible reading is a blind spot of its curvature model: the sampled Fisher rates a direction flat that the true past-task loss punishes further out, and new-task pressure escapes through it. The exposure is characteristic of feedforward control, whose gate can only be as good as its model of the past task; a feedback gate instead reads the damage as it arrives, and its own weakness lies in the conditioning of the signal it reads, as the three-task sequence shows (Section~\ref{subsec:disc_limits}). The cliffs and the plasticity debt are thus the two gates' defining and complementary costs, the feedforward gate cheap in plasticity but myopic, the feedback gate harder to blindside but slow, and this is why the two flip with momentum rather than sitting on a fixed frontier (Section~\ref{subsec:res_regimeflip}). Both gates carry to the ResNet at matched accuracy and much improved worst-case retention (Table~\ref{tab:res_cifar}), though the momentum cross there is confounded (Section~\ref{subsec:disc_limits}); it is the gate principle, not either particular method, that generalises.

\subsection{What the account changes in the literature}
\label{subsec:disc_literature}

The two-level split settles what reads as a disagreement between the two accounts of the gap that Section~\ref{subsec:stability_gap} set out. \citet{delange2023continualevaluationlifelonglearning} trace the drop to the gradient imbalance at the boundary; \citet{hess2024complementaryperspectivescontinuallearning} show the gap survives training on the exact joint loss, and \citet{kao2021naturalcontinuallearningsuccess} derive the geometry that makes it do so. Each account is right about one level and silent on the other. The imbalance is the acute driver of the overshoot: balancing it away is by far the largest single lever on depth in the driver ladder ($19.5 \rightarrow 5.7$ pp), and it leaves the arc untouched. The persistence under exact gradients is the arc: D4 reproduces it under exact, balanced, noiseless replay and measures what remains, $1.9$ pp of depth and an area of $3.2$, a controlled counterpart of the joint-loss observation and a quantification of the analytical mechanism. Read this way the two accounts stop competing. Most of the measured depth is the overshoot the imbalance account describes; most of what no repair of the update can remove is the geometry the trajectory account describes; and a mitigation aimed at either alone leaves the other level standing, which is the practical content of the $2\times2$.

The unified update of \citet{sodhani2022introductionlifelongsupervisedlearning} sorts replay methods by which weight they move and from what signal; the results add the axis that turned out to be decisive, where the modulation is applied. A-GEM and MEGA hold the new-task weight at one and correct per step from a live signal \citep{chaudhry2019efficientlifelonglearningagem, guo2020improvedschemesepisodicmemorybased}, which places them on the step side of the source/step divide, the side whose protection momentum unwinds in Section~\ref{subsec:res_directional}. The account therefore predicts that their boundary protection weakens under momentum as the directional gate's does, while source-attenuating schedules compose with it. These results do not test that prediction, but it is directly testable, and it matters because the original comparisons predate the continual-evaluation lens: end-of-task metrics would not have registered per-step protection eroding, any more than they register the gap itself.

Mode connectivity enters the account as an existence result and exits as a control result. \citet{garipov2018losssurfacesmodeconnectivity} and, for the continual setting, \citet{mirzadeh2020linearmodeconnectivitymultitask} establish post hoc that the relevant solutions share a low-loss corridor; the corridor is inspected after training and says nothing about whether an optimiser can be kept inside it while learning. C8 supplies the dynamic counterpart: an update rule whose iterate tracks the monotone reference path through the transition ($0.2$ pp of depth), turning the observation that a benign route exists into the statement that a gate can make training follow it. The same experiment sharpens the stability--plasticity dilemma \citep{Mermillod2013Stability-Plasticity} into a measurement. At fixed capacity and memory, the dilemma's price is the permanent $5.5$ pp offset to the joint level; the transient that continual evaluation reveals, up to $19.5$ pp at the boundary here, is not part of that price but a routing failure the right gate removes. The dilemma bounds where the model can settle; it does not oblige the model to swing below that level on the way.

\subsection{Limitations}
\label{subsec:disc_limits}

Several limits bound these claims. The first is statistical: with five seeds the design cannot reach a $0.05$ threshold (Section~\ref{subsec:metrics}), so the argument rests on effect sizes and seed-paired direction. The effects the story leans on, the up-to-eightfold depth reductions on CIFAR, the C7$\rightarrow$C8 variance collapse, the momentum re-inflation of the directional gate, are large relative to the seed spread and unanimous across seeds; the smaller ones, the adjacent-$\delta$ steps and the $\lambda_{\min}$ trade-offs, are read as directional consistency only, and one, the adaptive-versus-best-linear depth, does not reach even that ($p = 0.44$).

The second is the accuracy drift of the scalar gate at momentum $0$: lengthening the linear ramp lowers depth but costs several points of final accuracy (Table~\ref{tab:res_curriculum}), so the method's accuracy-neutrality is a property of the full curriculum-plus-momentum configuration, not of the ramp alone. The third is the directional gate's cliffs: the solver control establishes that they are real and not a numerical artefact, but the blind-spot mechanism we offer for them is a reading consistent with the data, not an established cause. The fourth is the CIFAR momentum cross, which we read as uninformative in both directions: at momentum $0$ the ResNet under-converges, both gates improve with momentum by similar margins, and the vanilla null sits on the same degraded baseline, so no cell of the deep-backbone cross cleanly confirms or refutes a momentum signature. All three momentum signatures therefore rest on the rot-MNIST results; a clean deep-backbone test would require a converged momentum-$0$ baseline we did not run. The fifth is the three-task group-norm sequence, where the adaptive ratio loses conditioning at a boundary between two near-identical corruptions and the schedule underperforms its control with high seed variance; we flag it as a reliability boundary of the gradient-ratio feedback signal, absent from the batch-norm variant, rather than a property of the gate principle; Appendix~\ref{app:cifar_gen_full} traces the failure step by step through the logged gradient-norm ratio.

Two scope limits are worth stating plainly. The full-buffer legs of both the driver block and the $\delta$ sweep store the entire past task and so break the limited-memory premise of continual learning; they are diagnostic probes that isolate the arc and split momentum's roles, and the practical configurations are the $1$k-buffer curriculum (C7$_M$) and the standard-buffer directional gate. And the directional gate's per-step Fisher-vector solve adds a fixed multiple of the backward cost per step, ${\sim}16\times$ the wall-clock and ${\sim}24\times$ the GPU energy of the curriculum on the ResNet-18 (Table~\ref{tab:res_cost}), an overhead the scalar gate's single gradient-norm ratio does not carry. The headline analysis also rests on a single-transition, two-task design chosen to isolate one boundary; the five-task rot-MNIST sequence (Appendix~\ref{app:longseq_full}) shows the regime picture surviving four consecutive boundaries, but a systematic study of longer sequences and of class-incremental rather than domain-incremental shifts is left open.

\subsection{Outlook}
\label{subsec:disc_outlook}

The taxonomy that organises the two methods also names what is missing from it. The gates studied here fill three of its four corners: scalar$\times$feedback, directional$\times$feedforward, and, as an ablation, the linear ramp at scalar$\times$feedforward. The remaining corner, a directional gate driven by a live damage signal rather than by a precomputed curvature model, is exactly the design the cliffs argue for: it would filter the push by direction, keeping the plasticity that makes the feedforward gate cheap, while reading the damage as it arrives, closing the blind spot that makes it myopic. Building and testing such a gate is the most direct consequence of these results, and the momentum results constrain it in advance: its directional attenuation must act at source or on the accumulated velocity, not on the instantaneous step, or it inherits the re-inflation of Section~\ref{subsec:res_directional}.

Three narrower directions follow from the failure modes and costs. First, the same repair applies to existing methods: any approach that preconditions the update with past-task curvature per step, the natural-gradient family included \citep{kao2021naturalcontinuallearningsuccess}, carries the re-inflation exposure, and moving its attenuation to the accumulated velocity is a directly testable fix. Second, the directional gate's compute cost is not fundamental: the per-step conjugate-gradient solve could be replaced by a diagonal or Kronecker-factored Fisher approximation, and whether the protection survives the coarser metric is open; on backbones much larger than the ResNet-18 here, such an approximation is a precondition for the directional corner rather than a refinement of it. Third, momentum's gap-closing role proved interchangeable with source variance reduction (C8), so the curriculum could be paired with other variance reducers; iterate averaging, which carries none of momentum's accumulator-driven overshoot, is the natural candidate.

The cliffs add a methodological rule. They arrive late and abruptly, after long stretches of exemplary retention, and end-of-task metrics would have recorded the $\delta{=}0.03$ runs as successes; only per-step measurement caught the failure. Continual evaluation \citep{delange2023continualevaluationlifelonglearning} is usually presented as the lens that reveals the stability gap. These results argue for it also as an acceptance test for any feedforward gate deployed in a system that serves predictions while it trains: a gate that acts on a model of the past task can fail precisely where that model is wrong, silently until it does.
